% GENERATED FILE. Edit canonical lessons, then run npm run generate:books.
% canonical-source-hash: fnv1a64:a73e9ca92fa617f4
% canonical-lessons: TE-C309-gattu, TE-C309-pasuvulu, TE-C309-stambham, TE-C309-kulayi, TE-C309-jenda, TE-R309-first-pass-words-from-chapters-305-306, TE-R309-second-pass-words-from-chapters-307-309

\chapter{Bund, Cattle, Pillar, Water tap, Flag}
\label{ch:te-bund-cattle-pillar-water-tap-flag}
\hlchaptermodality{309}

\begin{chapteropening}
\textbf{By the end of this chapter:} \emph{I can say a bund, cattle, a pillar, a water tap and a flag.}

Complete the last lesson of chapter 309: I can say a bund, cattle, a pillar, a water tap and a flag.
\end{chapteropening}

\section[\texorpdfstring{\te{గట్టు} (gaṭṭu)}{గట్టు (gaṭṭu)}]{\te{గట్టు} (gaṭṭu) — a bund, an embankment}

\label{lesson:TE-C309-gattu}



\begin{quote}
Before the new one: say the Telugu for a knot, then the Telugu for a row.
\end{quote}

\subsection*{You'll want to know: \te{గట్టు}}
\textbf{\te{గట్టు}} — \emph{gaṭṭu} — \textquotedblleft{}a bund, an embankment\textquotedblright{}.

\begin{grammarlens}[title={What you've built}]
One more word for this chapter.
\end{grammarlens}

\subsection*{Your turn}
\begin{itemize}
\raggedright
  \item \emph{Say it:} \emph{gaṭṭu}
  \item \emph{Say it:} \emph{gaṭṭu}, once more
  \item \emph{Say it:} say \emph{gaṭṭu}
  \item \emph{Recall:} say the Telugu for a knot, then the Telugu for a row, then say \emph{gaṭṭu} again
\end{itemize}

\begin{tcolorbox}[breakable,colback=teal!4,colframe=teal!35!black,title={Before you move on}]
What does \te{గట్టు} mean? (\textquotedblleft{}A bund, an embankment\textquotedblright{}.) Say it once more.
\end{tcolorbox}
\section[\texorpdfstring{\te{పశువులు} (paśuvulu)}{పశువులు (paśuvulu)}]{\te{పశువులు} (paśuvulu) — cattle, livestock}

\label{lesson:TE-C309-pasuvulu}



\begin{quote}
Before the new one: say the Telugu for a row, then the Telugu for a bund.
\end{quote}

\subsection*{You'll want to know: \te{పశువులు}}
\textbf{\te{పశువులు}} — \emph{paśuvulu} — \textquotedblleft{}cattle, livestock\textquotedblright{}.

\begin{grammarlens}[title={What you've built}]
One more word for this chapter.
\end{grammarlens}

\subsection*{Your turn}
\begin{itemize}
\raggedright
  \item \emph{Say it:} \emph{paśuvulu}
  \item \emph{Say it:} \emph{paśuvulu}, once more
  \item \emph{Say it:} say \emph{paśuvulu}
  \item \emph{Recall:} say the Telugu for a row, then the Telugu for a bund, then say \emph{paśuvulu} again
\end{itemize}

\begin{tcolorbox}[breakable,colback=teal!4,colframe=teal!35!black,title={Before you move on}]
What does \te{పశువులు} mean? (\textquotedblleft{}Cattle, livestock\textquotedblright{}.) Say it once more.
\end{tcolorbox}
\section[\texorpdfstring{\te{స్తంభం} (stambhaṁ)}{స్తంభం (stambhaṁ)}]{\te{స్తంభం} (stambhaṁ) — a pillar, a pole}

\label{lesson:TE-C309-stambham}



\begin{quote}
Before the new one: say the Telugu for a bund, then the Telugu for cattle.
\end{quote}

\subsection*{You'll want to know: \te{స్తంభం}}
\textbf{\te{స్తంభం}} — \emph{stambhaṁ} — \textquotedblleft{}a pillar, a pole\textquotedblright{}.

\begin{grammarlens}[title={What you've built}]
One more word for this chapter.
\end{grammarlens}

\subsection*{Your turn}
\begin{itemize}
\raggedright
  \item \emph{Say it:} \emph{stambhaṁ}
  \item \emph{Say it:} \emph{stambhaṁ}, once more
  \item \emph{Say it:} say \emph{stambhaṁ}
  \item \emph{Recall:} say the Telugu for a bund, then the Telugu for cattle, then say \emph{stambhaṁ} again
\end{itemize}

\begin{tcolorbox}[breakable,colback=teal!4,colframe=teal!35!black,title={Before you move on}]
What does \te{స్తంభం} mean? (\textquotedblleft{}A pillar, a pole\textquotedblright{}.) Say it once more.
\end{tcolorbox}
\section[\texorpdfstring{\te{కుళాయి} (kuḷāyi)}{కుళాయి (kuḷāyi)}]{\te{కుళాయి} (kuḷāyi) — a water tap, a faucet}

\label{lesson:TE-C309-kulayi}



\begin{quote}
Before the new one: say the Telugu for cattle, then the Telugu for a pillar.
\end{quote}

\subsection*{You'll want to know: \te{కుళాయి}}
\textbf{\te{కుళాయి}} — \emph{kuḷāyi} — \textquotedblleft{}a water tap, a faucet\textquotedblright{}.

\begin{grammarlens}[title={What you've built}]
One more word for this chapter.
\end{grammarlens}

\subsection*{Your turn}
\begin{itemize}
\raggedright
  \item \emph{Say it:} \emph{kuḷāyi}
  \item \emph{Say it:} \emph{kuḷāyi}, once more
  \item \emph{Say it:} say \emph{kuḷāyi}
  \item \emph{Recall:} say the Telugu for cattle, then the Telugu for a pillar, then say \emph{kuḷāyi} again
\end{itemize}

\begin{tcolorbox}[breakable,colback=teal!4,colframe=teal!35!black,title={Before you move on}]
What does \te{కుళాయి} mean? (\textquotedblleft{}A water tap, a faucet\textquotedblright{}.) Say it once more.
\end{tcolorbox}
\section[\texorpdfstring{\te{జెండా} (jeṇḍā)}{జెండా (jeṇḍā)}]{\te{జెండా} (jeṇḍā) — a flag}

\label{lesson:TE-C309-jenda}



\begin{quote}
Before the new one: say the Telugu for a pillar, then the Telugu for a water tap.
\end{quote}

\subsection*{You'll want to know: \te{జెండా}}
\textbf{\te{జెండా}} — \emph{jeṇḍā} — \textquotedblleft{}a flag\textquotedblright{}.

\begin{grammarlens}[title={What you've built}]
One more word for this chapter.
\end{grammarlens}

\subsection*{Your turn}
\begin{itemize}
\raggedright
  \item \emph{Say it:} \emph{jeṇḍā}
  \item \emph{Say it:} \emph{jeṇḍā}, once more
  \item \emph{Say it:} say \emph{jeṇḍā}
  \item \emph{Recall:} say the Telugu for a pillar, then the Telugu for a water tap, then say \emph{jeṇḍā} again
\end{itemize}

\begin{tcolorbox}[breakable,colback=teal!4,colframe=teal!35!black,title={Before you move on}]
What does \te{జెండా} mean? (\textquotedblleft{}A flag\textquotedblright{}.) Say it once more.
\end{tcolorbox}
\section[(dialogue)]{Words from chapters 305-306 — a first pass}

\label{lesson:TE-R309-first-pass-words-from-chapters-305-306}



\begin{quote}
Say the words for a water tap and a flag.
\end{quote}

\subsection*{Your turn}
Say these aloud:

\begin{itemize}
\raggedright
  \item an app, a keyboard, a laptop, a battery, a video
  \item steel, cement, leather, cotton, wool
\end{itemize}

\begin{tcolorbox}[breakable,colback=teal!4,colframe=teal!35!black,title={Before you move on}]
An app? (\textbf{\te{యాప్}}, \emph{yāp}.) A keyboard? (\textbf{\te{కీబోర్డు}}, \emph{kībōrḍu}.) A laptop? (\textbf{\te{ల్యాప్టాప్}}, \emph{lyāpṭāp}.) A battery? (\textbf{\te{బ్యాటరీ}}, \emph{byāṭarī}.) A video? (\textbf{\te{వీడియో}}, \emph{vīḍiyō}.) Steel? (\textbf{\te{ఉక్కు}}, \emph{ukku}.) Cement? (\textbf{\te{సిమెంటు}}, \emph{simeṇṭu}.) Leather? (\textbf{\te{తోలు}}, \emph{tōlu}.) Cotton? (\textbf{\te{పత్తి}}, \emph{patti}.) Wool? (\textbf{\te{ఉన్ని}}, \emph{unni}.)
\end{tcolorbox}
\section[(dialogue)]{Words from chapters 307-309 — a second pass}

\label{lesson:TE-R309-second-pass-words-from-chapters-307-309}



\begin{quote}
Say the words for a fridge and a mixer-grinder.
\end{quote}

\subsection*{Your turn}
Say these aloud:

\begin{itemize}
\raggedright
  \item a fridge, a mixer-grinder, dirty dishes, an incense stick, kumkum
  \item a pit, a piece, an edge, a knot, a row
  \item a bund, cattle, a pillar, a water tap, a flag
\end{itemize}

\begin{tcolorbox}[breakable,colback=teal!4,colframe=teal!35!black,title={Before you move on}]
A fridge? (\textbf{\te{ఫ్రిజ్జు}}, \emph{phrijju}.) A mixer-grinder? (\textbf{\te{మిక్సీ}}, \emph{miksī}.) Dirty dishes? (\textbf{\te{అంట్లు}}, \emph{aṇṭlu}.) An incense stick? (\textbf{\te{అగరుబత్తి}}, \emph{agarubatti}.) Kumkum? (\textbf{\te{కుంకుమ}}, \emph{kuṅkuma}.) A pit? (\textbf{\te{గుంట}}, \emph{guṇṭa}.) A piece? (\textbf{\te{ముక్క}}, \emph{mukka}.) An edge? (\textbf{\te{అంచు}}, \emph{añcu}.) A knot? (\textbf{\te{ముడి}}, \emph{muḍi}.) A row? (\textbf{\te{వరుస}}, \emph{varusa}.) A bund? (\textbf{\te{గట్టు}}, \emph{gaṭṭu}.) Cattle? (\textbf{\te{పశువులు}}, \emph{paśuvulu}.) A pillar? (\textbf{\te{స్తంభం}}, \emph{stambhaṁ}.) A water tap? (\textbf{\te{కుళాయి}}, \emph{kuḷāyi}.) A flag? (\textbf{\te{జెండా}}, \emph{jeṇḍā}.)
\end{tcolorbox}
